\begin{textsample}{POS Dim 3 – vem\_ed\_21 – Score 4.00 – vem\_ed\_21/t109.txt}  \label{ex:f3_pos_107}
Moderações Vozes emergentes, línguas e aprendizagem global Osvaldo Succi Junior fez a \textbf{moderação} de duas sessões de comunicação \textbf{oral}: “Leituras colaborativas e aprendizagem de línguas estrangeiras no projeto Literatandem”, por Vanessa Matiola (Unesp / Brasil) e “COIL en la enseñanza y el aprendizaje de idiomas extranjeros y segundas lenguas” (COIL no ensino e na aprendizagem de idiomas estrangeiros e segundas línguas), por Lilian Velasquez (DUOC UC / Colômbia) e Gabriel Cabezas (Universidad del Bío-Bío / Colômbia).
Ana Carolina Freschi fez a \textbf{moderação} do FlashTalk “Creative Global Learning at SHU: Focus on students’ engagement” (Aprendizagem Global Criativa na SHU: Foco no envolvimento dos alunos).
E apoiou Divinia Jithoo (Durban University of Technology / África do Sul) na \textbf{moderação} do \textbf{painel} “Emerging Voices” (Vozes Emergentes).
Estiveram presencialmente, no auditório principal do Centro de Convenções Rebouças, as vozes emergentes das painelistas Sandra Julieth Valencia Escobar (Universidad Católica De Manizales / Colômbia), Alia Gilbrecht (AN-Najah University / Palestina) e Bee Gan (Sheffield Hallam University / Reino Unido).
Daniel Otieno (Kenyatta University, Quênia) participou remotamente do \textbf{painel}.
Sandra Valencia destacou a importância de líderes de internacionalização nas instituições “sem os quais, o que ocorre são apenas iniciativas isoladas”.
E elencou os principais obstáculos à internacionalização sob a perspectiva latino-americana: barreira linguística, recursos econômicos, ferramentas tecnológicas e conectividade e perspectiva política.
Daniel Otieno ressaltou os princípios que devem guiar o design de projetos COIL: equidade, tomada de decisão democrática e translinguagem.
Alia Gilbrecht citou como principais desafios a captação de \textbf{fundos}, a conectividade e questões geopolíticas (como a guerra afeta a todos, física e emocionalmente).
Bee Gan sugere redesenhar os cursos para incluir a internacionalização no currículo.
Do ponto de vista estratégico, recomendou: “olhe para a missão da universidade, seus KPIs (Key Performance Indicators — indicadores-chave de desempenho) e mostre como os projetos COIL ajudam nisso.
Dessa forma, os Intercâmbios Virtuais ganham adesão institucional”.

% matched lemmas: fundo, moderação, oral, painel
\end{textsample}
